%قبل از شروع مسئله ترجیحاً چیزی ننویسید. اما اگر خواستید بنویسید، در قالب لیستی (فقرات و شمارش) بنویسید.
%برای اضافه کردن نمره، بعدی از شروع‌مسئله، [x] را اضافه کنید و به جای x  نمره مورد نظر را قرار دهید. برای زیرمسئله‌ها هم به همین صورت است.
%\شروع‌مسئله[10]
\شروع‌مسئله [25]
% صورت اصلی مسئله را در این‌جا بنویسید.
\متن‌سیاه{دگرگونی}

فرض کنید که $s$ یک رشته با طول $n$ باشد. برای اعداد صحیح $i$ و $j$ که 
$0 \leq i < j < n$
تابع $f$ را به صورت زیر تعریف می‌کنیم:
$$f(s, i, j) = s[i+1:j-1] + r(s[j:n-1]) + r(s[0:i]) $$
که در اینجا منظور از 
$s[p:q]$
زیررشته‌ای از $s$ است که از $p$ شروع و به $q$ ختم می‌شود. علامت $+$ به معنی ادغام دو رشته است و همچنین تابع $r$ معکوس رشته ورودی را برمی‌گرداند. (اگر 
$j=i+1$
زیر رشته 
$s[i+1:j-1]$
تهی در نظر گرفته می‌شود)

با این تعاریف، به شما دو رشته $a$ و $b$ به عنوان ورودی داده می‌شود و به عنوان خروجی باید $i$ و $j$ای را به عنوان خروجی چاپ کنید که  
$f(a, i, j) = b$
به طوری که ‌$i$ بزرگترین مقدار ممکن باشد و اگر برای این $i$ چند مقدار $j$ وجود داشت، کم‌ترین 
مقدار $j$ را انتخاب کنید.

\متن‌سیاه{ورودی:}\\
در دو خط اول ورودی، دو رشتۀ ناتهی $a$ و $b$ با طول کم‌تر از $10^6$ به عنوان ورودی داده می‌شوند. (رشته‌ها می‌توانند شامل هر کاراکتری با کد \متن‌لاتین{ASCII} بین ۳۲ و ۱۲۶ باشند) 

\متن‌سیاه{خروجی:}\\
دو عدد $i$ و $j$ به عنوان خروجی چاپ می‌شوند. اگر به ازای ورودی جوابی وجود نداشت، عبارت 
\متن‌تایپ{-1 -1}
چاپ می‌شود.\\

\متن‌سیاه{ورودی و خروجی نمونه:}\\
\begin{latin}
	\begin{tabularx}{\textwidth}{|X | X|}
		\hline
		\textbf{input} &
		\textbf{output}\\
		\hline
		\texttt{cbaaaa}
		&
		\texttt{4 5}\\
		\texttt{aaaabc}&\\
		\hline
		
	\end{tabularx}
\end{latin}

\begin{latin}
	\begin{tabularx}{\textwidth}{|X | X|}
		\hline
		\textbf{input} &
		\textbf{output}\\
		\hline
		\texttt{123342}
		&
		\texttt{-1 -1}\\
		\texttt{3324212}&\\
		\hline
		
	\end{tabularx}
\end{latin}


\begin{latin}
	\begin{tabularx}{\textwidth}{|X | X|}
		\hline
		\textbf{input} &
		\textbf{output}\\
		\hline
		\texttt{abc}
		&
		\texttt{1 2}\\
		\texttt{cba}&\\
		\hline
		
	\end{tabularx}
\end{latin}

%\شروع{زیرمسائل}
%\زیرمسئله[0] % مسئلۀ ۱آ
%\زیرمسئله % مسئلۀ ۱ب
%\زیرمسئله % مسئلۀ ۱پ
%\زیرمسئله % مسئلۀ ۱ت
%\پایان{زیرمسائل}

\solution{
رشتۀ $c$ را معکوس رشتۀ $b$ در نظر بگیرید. در این صورت باید $i$ و $j$ای را بیابیم که 
$$ c = a[0 : i] + a[j : n-1] + r(a[i+1 : j-1])$$

برای این کار بیشترین $i$ای که
$ c[0 : i] = a[0 : i]$
را پیدا می کنیم. سپس به ازای هر $i$ بین 0 تا این مقدار بیشینه، اگر بیشترین مقدار $k$ را پیدا کنیم که
$ c[i+1 : k]  = a[n-(k-i) : n-1]$
می توانیم چک کنیم که خواستۀ مسئله برقرار است یا نه (این که به ازای $i$ فقط بیشترین $k$ کافیست به‌خاطر کمینه بودن $j$ در حالتی که چند جواب برای یک $i$ داشتیم است).

حال برای پیدا کردن چنین $k$ بیشینه‌ای برای هر$i$، به صورت برعکس عمل می‌کنیم و به ازای هر $k$ کوچک‌ترین $i$ای را می‌یابیم که
$ c[i+1 : k]  = a[n-(k-i) : n-1]$
برقرار باشد. برای این کار کافیست روی $i$،
جست‌وجوی دودویی 
انجام دهیم (چون اگر معادله بالا به ازای یک $i$ برقرار باشد به ازای مقادیر بیشتر از $i$ نیز صادق است، بنابراین می‌توانیم از جست‌وجوی دودویی استفاده کنیم.)

چک کردن این که دو زیر رشته از $a$ و $c$ آیا با هم برابرند یا نه هم را با نگه داشتن درهم‌سازی هر  پیشوند از 0 تا $i$ از رشته‌های $a$ و $c$ می‌توان انجام داد. چون با داشتن این مقادیر درهم‌سازی از $i$ تا $j$ رشته را می توانیم در $O(1)$ بدست بیاوریم.

بنابراین پیچیدگی زمانی هم از $O(nlog(n))$ می باشد.

کد آقای رضا یاربخش با نام \متن‌لاتین{Degarguni\_sol.py} به عنوان نمونه در پوشه مربوطه قرار داده شده. 

}




%%%%%%%%%%%%%%%%%%%%%%%%%%%%%%%%%%%%%%%%%%%%%%%%%%%%%
% برخی دستورات مفید برای تمرین‌ها. %
%%%%%%%%%%%%%%%%%%%%%%%%%%%%%%%%%%%%%%%%%%%%%%%%%%%%%

% فرمول‌های نمونه:
%$f(x) = \Theta(x)$
%$T(x, y) \leq \log(x) + 2^y + \binom{2n}{n}$
% {\tt code\_function}


% لیست موارد بدون شماره‌گذاری:
%\شروع{فقرات}
%    \فقره مورد اول 
%        \شروع{فقرات}
%            \فقره مورد اول داخلی 
%                \شروع{فقرات}
%                    \فقره اولین مورد خیلی داخلی 
%                    \فقره دومین مورد خیلی داخلی 
%                \پایان{فقرات}
%            \فقره مورد دوم داخلی 
%        \پایان{فقرات}
%    \فقره مورد دوم 
%\پایان{فقرات}

% لیست موارد شماره‌گذاری شده:
%\شروع{شمارش}
%    \فقره یک مورد 
%    \فقره یک مورد دیگر 
%    \فقره باز هم یک مورد دیگر
%\پایان{شمارش}

% برای نوشتن یک فرمول چند خطی که یک شماره بخوره. برای حالت‌های دیگر به سایت زیر مراجعه کنید:
% https://www.overleaf.com/learn/latex/Aligning_equations_with_amsmath
%\begin{align} 
%    \begin{split}
%        (x+y)^3 &= (x+y)^2(x+y) \\
%                &= (x^2+2xy+y^2)(x+y) \\
%                &= (x^3+2x^2y+xy^2) + (x^2y+2xy^2+y^3) \\
%                &= x^3+3x^2y+3xy^2+y^3
%    \end{split}                                 
%\end{align}

% نوشتن ماتریس و آرایه دوبعدی
%\begin{align}
%    A = 
%    \begin{bmatrix}
%        A_{11} & A_{21} \\
%        A_{21} & A_{22}
%    \end{bmatrix}
%\end{align}

% اضافه کردن تصویر با شرح. برای اضافه کردن تصویر بدون شرح، فقط دستور درج‌تصویر را استفاده کنید.
%\شروع{شکل}[h!]
%\تنظیم‌ازوسط
%  \درج‌تصویر[پهنا=0.5\پهنای‌متن]{img.jpg}
%\شرح{شرح تصویر}
%\برچسب{تصویر:۱}
%\پایان{شکل}
%\پایان{زیرمسائل}

%کد زیر را ببینید.
%%\lstinputlisting{source_filename.py}
%\begin{latin}
%\begin{lstlisting}
%# Import the modules
%import sys
%import random
%
%ans = True
%
%while ans:
%question = raw_input("Ask the magic 8 ball a question: (press enter to quit) ")
%
%answers = random.randint(1,8)
%
%if question == "":
%sys.exit()
%
%elif answers == 1:
%print "It is certain"
%
%elif answers == 2:
%print "Outlook good"
%
%elif answers == 3:
%print "You may rely on it"
%
%elif answers == 4:
%print "Ask again later"
%
%elif answers == 5:
%print "Concentrate and ask again"
%
%elif answers == 6:
%print "Reply hazy, try again"
%
%elif answers == 7:
%print "My reply is no"
%
%elif answers == 8:
%print "My sources say no"
%\end{lstlisting}
%\end{latin}
